% ===================================================================
%  TAULUKKO N:o 4 — Kypsä ikä ja vanhuus.
%  Traduction 2026 d'apres texto/io, controlee sur texto/fr et
%  texto/sv. Consigne : voir texto/fi/10-taulukko-01.tex.
% ===================================================================

%%K t04-apar-1 apar
\VUsaut{4.0mm}
\VUtitre{15.0pt}{60}{TAULUKKO N:o 4}
\VUsaut{1.6mm}
\VUtitre{11.4pt}{0}{Kypsä ikä ja vanhuus.}
\VUsaut{3.2mm}

%%K t04-tit-1 sub
\VUcentre{11.4pt}{0}{\textit{Ensimmäinen kohtaus.}}
\VUsaut{1.6mm}
\VUtitre{11.4pt}{0}{Hääateria.}
\VUsaut{2.6mm}

%%K t04-tit-2 sub
\VUpk{10.2pt}{40}{Syksy.}
\VUsaut{3.0mm}

%%K t04-c1-01-1 p
1. --- Nuoret ovat kasvaneet aikuisiksi. Yksi heistä, Juhani, menee naimisiin.
Olemme läsnä \VUgras{hääaterialla}, joka kokoaa puolisoiden perheet saman pöydän
ympärille.

%%K t04-c1-02-1 p
2. --- Olemme \VUgras{syksyssä}, \VUgras{syyskuussa}. \VUgras{Lokakuun} ja
\VUgras{marraskuun} kylmä tuuli ei ole vielä riistänyt maaseudulta sen vihreää
lehvistöä. Illan ilma on haalea ja tuoksuva; siksi ateria pidetään puutarhan
puoleisen \VUgras{kuistin} \textsuperscript{(8)} alla.

%%K t04-c1-03-1 p
3. --- Kaukana, \VUgras{kukkulalla} \textsuperscript{(3)}, on Juhanin vanhempien
koti, tyylikäs \VUgras{linna} \textsuperscript{(1)}, joka on kätkeytynyt
vehreyteen; kauniit viinitarhat, jotka antavat erinomaista viiniä, peittävät
kukkulan rinteen. Viininviljelijä viljelee ja hoitaa niitä.

%%K t04-c1-04-1 p
4. --- Viinitarhojen yllä lentää parvi \VUgras{peltopyitä} \textsuperscript{(2)},
jotka \VUgras{riistanvartijan} \textsuperscript{(5)} lähestyminen on
säikäyttänyt. Tämä tutkii, \VUgras{kivääri} kainalossa ja
\VUgras{metsästyslaukku olalla}, \VUgras{pensaikkoa} \textsuperscript{(4)}
\VUgras{koiransa} \textsuperscript{(6)} kanssa. Hän valvoo tilaa ja estää
salametsästäjiä hävittämästä riistaa.

%%K t04-c1-05-1 p
5. --- \VUgras{Kuistin} ulkopuolella kiipeää \VUgras{säleikköviiniköynnös},
josta riippuu tuuheita \VUgras{rypäleterttuja} \textsuperscript{(7)}.

%%K t04-c1-06-1 p
6. --- Kauniit syyskukat, jotka koristavat \VUgras{pöytää}
\textsuperscript{(16)}, ovat \VUgras{krysanteemeja} \textsuperscript{(9)};
morsiamen \VUgras{isä} \textsuperscript{(10)} on pannut yhden niistä frakkinsa
napinläpeen.

%%K t04-c1-07-1 p
7. --- Ateria lähenee loppuaan. \VUgras{Palvelija} \textsuperscript{(11)} alkaa
kantaa sisään jälkiruokavateja ja poistaa pian kattauksen eri osat,
\VUgras{lusikat} \textsuperscript{(14)}, \VUgras{haarukat}
\textsuperscript{(12)}, \VUgras{veitset} \textsuperscript{(13)}, \VUgras{suuret
vadit} \textsuperscript{(15)}, joita ei enää tarvita.

%%K t04-tit-3 sub
\VUpk{10.2pt}{40}{Suku.}
\VUsaut{3.0mm}

%%K t04-c1-08-1 p
8. --- Kaikki näyttävät iloisilta, ja hymy, jolla sulhasen \VUgras{äiti}
\textsuperscript{(17)} katselee lapsiaan, todistaa hänen onnestaan.

%%K t04-c1-09-1 p
9. --- Tuo avioliitto yhdistää kaksi seudun rikasta perhettä. Juhani,
\VUgras{aviomies} \textsuperscript{(18)}, on suuren maanomistajan
\VUgras{poika} ja hänestä tulee taitavan teollisuusmiehen \VUgras{vävy}.

%%K t04-c1-10-1 p
10. --- \VUgras{Puolisot} ovat nuoria ja kuitenkin he olivat jo kauan
\VUgras{kihloissa}. \VUgras{Nuori nainen} \textsuperscript{(19)}, Martta, on
hurmaava \VUgras{valkoisessa puvussaan}, jota appelsiininkukat koristavat.

%%K t04-c1-11-1 p
11. --- Hänen \VUgras{appensa} \textsuperscript{(20)} on hyvin onnellinen
saadessaan niin rakastettavan \VUgras{miniän}, ja Juhanin \VUgras{anoppi}
\textsuperscript{(21)} hymyilee nähdessään \VUgras{tyttärensä} noin kauniina.

%%K t04-c1-12-1 p
12. --- Pöydän päässä istuvat muut \VUgras{vieraat} \textsuperscript{(22)},
\VUgras{sukulaiset, ystävät, serkut, serkut}. \VUgras{Hovimestari}
\textsuperscript{(23)}, lautasliina kainalossaan, palvelee heitä ahkerasti.

%%K t04-tit-4 sub
\VUpk{10.2pt}{40}{Ystävät.}
\VUsaut{3.0mm}

%%K t04-c1-13-1 p
13. --- Pöydän oikealla puolella, tässä on \VUgras{neiti}
\textsuperscript{(24)}, nuoren morsiamen \VUgras{ystävätär}, sitten hänen
\VUgras{veljensä}, joka on tämän \VUgras{sulhaspoika} \textsuperscript{(25)}.
Hän juttelee \VUgras{morsiusneitonsa} \textsuperscript{(26)} kanssa, joka on
sulhasen sisar ja josta tuon avioliiton kautta tulee Martan \VUgras{käly}. Hän
on hurmaava nuori tyttö. Hänellä on kauniita \VUgras{koruja, helminauha} ja
\VUgras{rannerengas} kultaa.

%%K t04-c1-14-1 p
14. --- Heidän vieressään herra, joka nousee seisomaan ja kohottaa lasinsa, on
perheen \VUgras{ystävä} \textsuperscript{(27)}. Hän on yksi \VUgras{sulhasen
todistajista}. Hän \VUgras{esittää maljan}, juo nuoren avioparin
\VUgras{terveydeksi} ja toivottaa heille pitkää onnea. Hän on pukeutunut
juhla-asuun, frakkiin ja \VUgras{lakkakenkiin} \textsuperscript{(28)}. Hän on
kavaljeeri \VUgras{isoäidille} \textsuperscript{(29)}, joka on \VUgras{leski}.

%%K t04-c1-15-1 p
15. --- Nuori mies \VUgras{frakissa} \textsuperscript{(31)}, jolla on ohuet
mustat viikset, on Juhanin \VUgras{lanko} \textsuperscript{(30)}. Hän on
huomattava insinööri. Hän istuu hyvin tyylikkään \VUgras{nuoren rouvan}
\textsuperscript{(33)} vieressä, joka on Martan \VUgras{vanhempi sisar} ja sen
sievän pikku tytön äiti, joka tulee, käsi serkkunsa ympärillä, tarjoamaan kukkaa
ja \VUgras{toivottamaan} hänelle \VUgras{hyvää iltaa}. Tuolla rouvalla on hyvin
kauniit \VUgras{korvarenkaat} ja tyylikäs \VUgras{päähine}.

%%K t04-c1-16-1 p
16. --- Pikku serkku, niin omalaatuinen \VUgras{puserossaan}
\textsuperscript{(36)} ja pussimaisissa \VUgras{housuissaan}
\textsuperscript{(37)}, on hyvin säälittävä. Hän on \VUgras{orpo}
\textsuperscript{(35)}. Hyvin nuorena hän menetti isänsä ja äitinsä. Hänen
setänsä on hänen \VUgras{holhoojansa} \textsuperscript{(38)}, ja tämä on pysynyt
naimattomana kasvattaakseen tuon poloisen lapsen. Hän on hyväsydäminen mies. Hän
on hieman kalju ja hyvin likinäköinen; hän käyttää \VUgras{nenälaseja}.

%%K t04-tit-5 sub
\VUsaut{2.4mm}
\VUpk{10.2pt}{40}{Jälkiruoka.}
\VUsaut{3.0mm}

%%K t04-c1-17-1 p
17. --- Vieraiden takana, pöydällä, joka on peitetty \VUgras{pöytäliinalla}, on
\VUgras{jälkiruoka} \textsuperscript{(39)} katettuna. Suuressa maljassa, tässä on
hedelmiä, \VUgras{persikoita} \textsuperscript{(40)}, \VUgras{päärynöitä}
\textsuperscript{(41)}, \VUgras{omenoita} \textsuperscript{(42)},
\VUgras{aprikooseja} \textsuperscript{(43)} ja \VUgras{viinirypäleitä}
\textsuperscript{(44)}. Aivan lähellä, \VUgras{ananaksia}
\textsuperscript{(45)}, kaunis \VUgras{kakku} \textsuperscript{(46)},
\VUgras{likööripulloja} \textsuperscript{(48)} ja \VUgras{samppanjaa}
\textsuperscript{(49)} sekä myös pinoja puhtaita \VUgras{lautasia}
\textsuperscript{(47)}.

%%K t04-c1-18-1 p
18. --- \VUgras{Kellarimestari} \textsuperscript{(50)} vetää auki vanhan viinin
\VUgras{pullon} \textsuperscript{(52)} \VUgras{korkkiruuvilla}
\textsuperscript{(51)}. \VUgras{Korkki} \textsuperscript{(53)} näyttää hyvin
vaikealta vetää ulos. Tuolla kellarimestarilla on \VUgras{lautasliina}
\textsuperscript{(54)} kainalossaan.

%%K t04-c1-19-1 p
19. --- Aterian jälkeen herrat menevät tupakkahuoneeseen, ja naiset menevät
laittautumaan valmiiksi tanssiaisiin.

%%K t04-tit-6 sub
\VUsaut{5.0mm}
\VUcentre{11.4pt}{0}{\textit{Toinen kohtaus.}}
\VUsaut{1.6mm}
\VUpk{11.4pt}{40}{Isoisän syntymäpäivä.}
\VUsaut{2.6mm}

%%K t04-tit-7 sub
\VUpk{10.2pt}{40}{Vierailu.}
\VUsaut{3.0mm}

%%K t04-c2-01-1 p
1. --- Kymmenen vuotta on kulunut, Juhanin vanhemmat ovat tulleet vanhoiksi ja
rakastavat kovasti kolmea lapsenlastaan, kahta \VUgras{pojanpoikaa}
\textsuperscript{(2)} ja yhtä \VUgras{pojantytärtä} \textsuperscript{(3)}. Koska
tänään on \VUgras{isoisän syntymäpäivä}, lapset tulevat toivottamaan heille
hyvää nimipäivää.

%%K t04-c2-02-1 p
2. --- Pikku Pietari on vielä hyvin nuori, hän on vasta kolmivuotias. Hän näkee
\VUgras{kissan} \textsuperscript{(1)} lähestyvän. Hän tahtoo silittää sen kaunista
silkinpehmeää \VUgras{turkkia} ja pitkää \VUgras{häntää}, hän ei ajattele, että
sulavien \VUgras{käpälien} alla piilee ilkeitä teräviä kynsiä, jotka ehkä
raapivat häntä.

%%K t04-c2-03-1 p
3. --- Hänen sisarensa Gabriella, joka antaa hänelle kätensä, on paljon
vakavampi, ja Marcel suurine \VUgras{takkeineen} \textsuperscript{(4)} ja
\VUgras{nahkavöineen} \textsuperscript{(5)} näyttää hieman miehekkäämmältä,
lyhyistä housuistaan huolimatta, jotka antavat nähdä hänen \VUgras{sukkansa}
\textsuperscript{(6)}.

%%K t04-c2-04-1 p
4. --- Heidän isänsä on pannut päälleen paksun \VUgras{talvitakkinsa}
\textsuperscript{(7)}, jossa on \VUgras{turkiskaulus} \textsuperscript{(8)}, ja
\VUgras{taskussaan} \textsuperscript{(9)} hänellä on kaulahuivi. Hänen
vaimollaan on huopahattunsa, jota koristavat \VUgras{höyhenet}
\textsuperscript{(10)}, \VUgras{jakkunsa} \textsuperscript{(11)},
\VUgras{turkisboansa} \textsuperscript{(12)} ja \VUgras{puuhkansa}
\textsuperscript{(13)}.

%%K t04-c2-05-1 p
5. --- On hyvin kylmä, sillä talvi vallitsee. \VUgras{Lunta}
\textsuperscript{(19)} on satanut koko päivän ja talojen katot ovat aivan
valkoiset. Yhdessä \VUgras{yön} kanssa kylmyys käy vieläkin purevammaksi.
Taivaalla loistavat \VUgras{kuu} \textsuperscript{(17)} ja \VUgras{tähdet}
\textsuperscript{(18)} ja lävistävät \VUgras{pimeyden}.

%%K t04-tit-8 sub
\VUsaut{4.2mm}
\VUpk{10.2pt}{40}{Sairaus.}
\VUsaut{3.0mm}

%%K t04-c2-06-1 p
6. --- Poloinen \VUgras{sokea} \textsuperscript{(20)}, joka kulkee torin yli
\VUgras{apteekin} \textsuperscript{(16)} edessä \VUgras{villakoiransa}
\textsuperscript{(21)} taluttamana, on hyvin onneton. Justiina,
\VUgras{palvelustyttö} \textsuperscript{(14)}, kantaa keittiöstä
\VUgras{kupillisen} \textsuperscript{(15)} hyvin kuumaa \VUgras{yrttiteetä},
joka on runsaasti sokeroitu.

%%K t04-c2-07-1 p
7. --- \VUgras{Isoäiti} \textsuperscript{(23)}, joka haluaa etsiä
\VUgras{lornettinsa} \textsuperscript{(26)} pöydältä lukeakseen
\VUgras{reseptin} \textsuperscript{(27)}, jonka \VUgras{lääkäri}
\textsuperscript{(25)} on juuri kirjoittanut, nousee nojatuolistaan ja nojaa
\VUgras{keppiinsä} \textsuperscript{(24)}.

%%K t04-c2-08-1 p
8. --- Usein hän kärsii \VUgras{huonovointisuudesta, huimauksesta,
vilustumisista, hammassärystä}, hänen jalkansa ovat heikot eivätkä hänen
silmänsä ole enää kovin hyvät. Siitä huolimatta hän laskee mielellään leikkiä
vanhan \VUgras{lääkärinsä} kanssa, joka aina tahtoo määrätä hänelle
\VUgras{lääkeliuosta} \textsuperscript{(28)}. Tänään hän on hyvin iloinen,
koska hänellä on nuori perheensä ympärillään.

%%K t04-c2-09-1 p
9. --- Hyvin suljetussa huoneessa on miellyttävää. Marmoroidussa
\VUgras{takassa} \textsuperscript{(29)} leimuaa \VUgras{komea tuli}
\textsuperscript{(30)}, ja tunnemme itsemme onnellisiksi \VUgras{lampun}
\textsuperscript{(31)} lempeässä \VUgras{valossa}, joka on juuri sytytetty.

%%K t04-tit-9 sub
\VUsaut{4.2mm}
\VUpk{10.2pt}{40}{Isoisä.}
\VUsaut{3.0mm}

%%K t04-c2-10-1 p
10. --- Jopa \VUgras{isoisä} \textsuperscript{(33)} on unohtanut, että hän oli
sairas, että hänen \VUgras{reumatisminsa} kahlitsee hänet
\VUgras{nojatuoliinsa} \textsuperscript{(34)} ja että hänellä vielä eilen oli
\VUgras{kuumetta ja pyörtymisiä}. Hän ei enää muista, että hän viime talvena
kärsi \VUgras{keuhkokuumeesta} ja että käheä ja tuskallinen \VUgras{yskä} sai
hänet kärsimään kovasti.

%%K t04-c2-10-2 p
Ja kuitenkin hänet parannettiin hyvin vaivoin. Nyt hän voi hieman paremmin.
Mutta hänen jalkaansa, jota hän tukee \VUgras{tyynyllä} \textsuperscript{(38)},
joka on asetettu pienelle \VUgras{jakkaralle} \textsuperscript{(39)}, koskee
myös.

%%K t04-c2-11-1 p
11. --- Kun hänet on huolellisesti kääritty lämpimään
\VUgras{yönuttuunsa} \textsuperscript{(35)}, jossa on paksut helmiäiset
\VUgras{napit} \textsuperscript{(36)}, ja kun hänellä on vieressään, lukemassa
hänelle lehteä, pieni \VUgras{orpotyttö} \textsuperscript{(40)}, joka on puettu
valkoiseen \VUgras{myssyyn}, siniseen \VUgras{pukuun} \textsuperscript{(42)} ja
\VUgras{viittaan} \textsuperscript{(41)}, hän ei valita.

%%K t04-c2-12-1 p
12. --- Joka vuosi, kun \VUgras{kalenteri} \textsuperscript{(43)} jälleen
näyttää \VUgras{päivämäärän} ensimmäinen \VUgras{joulukuuta}, hän
odottaa kärsimättömänä \VUgras{lapsenlapsiaan}, sillä hän tietää, etteivät he
unohda tuota tärkeää \VUgras{päivämäärää}.

%%K t04-c2-13-1 p
13. --- Hän muistaa liikuttuneena sen hyvin kaukaisen ajan, jolloin hän
henkilökohtaisesti tuli toivottamaan hyvää nimipäivää sille kunniakkaalle
keisarilliselle \VUgras{marsalkalle}, jonka \VUgras{muotokuva}
\textsuperscript{(44)} riippuu seinällä ja joka on hänen lastensa
\VUgras{isoisän isä}.
\VUsaut{6.0mm}
\VUfilet{22mm}
